\NormalFont\ShortTitle{Efésios}
{\MT Carta de Paulo aos Efésios

\par }\ChapOne{1}{\PP \VerseOne{1}Paulo, apóstolo de Jesus Cristo pela vontade de Deus, aos santos que estão em Éfeso, e fiéis em Cristo Jesus.
\VS{2}Graça e paz, de Deus nosso Pai, e do Senhor Jesus Cristo,
{\ADD{sejam}} convosco.
\VS{3}Bendito seja o Deus e Pai do nosso Senhor Jesus Cristo, que nos abençoou com todas as bênçãos espirituis nos lugares celestiais em Cristo.
\VS{4}Assim como nos escolheu nele antes da fundação do mundo, para que fôssemos santos e irrepreensíveis diante dele em amor.
\VS{5}E nos predestinou como filhos adotados por meio de Jesus Cristo para si mesmo, conforme o bom prazer de sua vontade;
\VS{6}para louvor da glória de sua graça, pela qual nos fez aceitáveis no Amado.
\VS{7}Nele temos a libertação
\FTNT{*}{{\FR 1:7 }
Ou: redenção, isto é, a libertação por meio do pagamento de um resgate
} pelo seu sangue, o perdão dos pecados, conforme as riquezas da sua graça,
\VS{8}que ele fez transbordar em nós em toda sabedoria e prudência.
\VS{9}E nos revelou o mistério da sua vontade, conforme o seu bom prazer, que propôs nele,
\FTNT{†}{{\FR 1:9 }
provavelmente em Cristo
}
\VS{10}para a administração do cumprimento dos tempos,
{\ADD{isto é,}} voltar a reunir todas as coisas em Cristo, tanto as que
{\ADD{estão}} nos céus, quanto as que
{\ADD{estão}} na terra.
\VS{11}Nele também fomos feitos herança, havendo sido predestinados conforme o propósito daquele que faz todas as coisas segundo o conselho da sua vontade,
\VS{12}para que fôssemos para louvor da sua glória, nós os que primeiro esperamos em Cristo.
\VS{13}E também vós, depois que ouvistes a palavra da verdade, o evangelho da vossa salvação, e nele crestes, nele fostes selados com o Espírito Santo da promessa,
\VS{14}que é a garantia da nossa herança, até a libertação
\FTNT{‡}{{\FR 1:14 }
Ou: redenção
} da propriedade
{\ADD{de Deus}} , para louvor da sua glória.
\VS{15}Por isso eu, desde que ouvi da fé que há entre vós no Senhor Jesus, e do amor a todos os santos,
\VS{16}não paro de agradecer
{\ADD{a Deus}} por vós, lembrando-me de vós em minhas orações.
\VS{17}{\ADD{Oro}} para que o Deus de nosso Senhor Jesus Cristo, o Pai glorioso
\FTNT{§}{{\FR 1:17 }
Lit. Pai da glória
} , vos dê o Espírito de sabedoria e de revelação no conhecimento dele;
\VS{18}e que sejam iluminados os olhos do vosso entendimento, para que conheçais qual é a esperança para a qual ele chamou, quais
{\ADD{são}} as riquezas da gloriosa herança
\FTNT{*}{{\FR 1:18 }
Lit. glória da herança
} dele nos santos,
\VS{19}e qual é a superabundante grandeza do seu poder para nós, que cremos, conforme a operação da sua poderosa força,
\VS{20}a qual ele operou em Cristo, ressuscitando-o dos mortos; e o colocou à direita dele nos céus,
\VS{21}acima de todo governo, autoridade, poder, domínio, e de todo nome que se nomeia, não só nesta era, mas também na futura.
\VS{22}Ele também sujeitou todas as coisas debaixo dos seus pés, e o constituiu por cabeça sobre todas as coisas para a Igreja,
\VS{23}que é o seu corpo, a plenitude daquele que enche tudo em todas as coisas.

\par }\Chap{2}{\PP \VerseOne{1}E vós estáveis mortos em ofensas e pecados,
\VS{2}nos quais antes andastes conforme o proceder deste mundo, conforme o príncipe do poder do ar, do espírito que agora opera nos filhos da desobediência.
\VS{3}Entre esses também todos nós antes andávamos nos desejos da nossa carne, fazendo a vontade da carne e dos pensamentos; e éramos por natureza filhos da ira, como também os outros.
\VS{4}Mas Deus, que é rico em misericórdia, pelo seu muito amor com que nos amou,
\VS{5}estando nós ainda mortos em
{\ADD{nossas}} ofensas, nos deu vida juntamente com Cristo (pela graça sois salvos),
\VS{6}{\ADD{nos}} ressuscitou, juntamente
{\ADD{com ele}} , e
{\ADD{nos}} fez sentar nos lugares celestiais, em Cristo Jesus;
\VS{7}para mostrar nos tempos futuros as abundantes riquezas da sua graça, pela
{\ADD{sua}} bondade conosco em Cristo Jesus.
\VS{8}Porque pela graça sois salvos, por meio da fé; e isto não
{\ADD{vem}} de vós;
{\ADD{é}} dom de Deus;
\VS{9}não por obras, para que ninguém tenha orgulho de si mesmo.
\VS{10}Pois nós fomos feitos por ele, criados em Cristo Jesus para as boas obras, que Deus preparou para que nelas andássemos.
\VS{11}Portanto, lembrai-vos de que vós, antes, éreis gentios na carne, e chamados de não-circuncidados pelos que se chamam participantes da circuncisão na carne, feita por mãos humanas;
\VS{12}que naquele tempo estáveis sem Cristo, excluídos da comunidade de Israel, e estranhos aos pactos da promessa, não tendo esperança, e sem Deus no mundo.
\VS{13}Mas agora em Cristo Jesus, vós que antes estáveis longe, chegastes para perto pelo sangue de Cristo,
\VS{14}pois ele é a nossa paz. Dos dois
{\ADD{povos}} ele fez um, e derrubou do meio o muro da separação.
\VS{15}Na sua carne ele desfez a inimizade da Lei dos mandamentos
{\ADD{que consistia}} em ordenanças, para criar em si mesmo os dois em um novo homem, fazendo a paz;
\VS{16}e pela cruz reconciliar com Deus a ambos num só corpo, matando nela a inimizade.
\VS{17}Ele veio, e proclamou o evangelho da paz a vós que
{\ADD{estáveis}} longe, e aos que estavam perto.
\VS{18}Pois por meio dele ambos temos acesso ao Pai por um
{\ADD{mesmo}} Espírito.
\VS{19}Portanto, já não sois estrangeiros, nem forasteiros, mas sim, concidadãos dos santos, e membros da família de Deus,
\VS{20}edificados sobre o fundamento dos apóstolos e dos profetas, do qual Jesus Cristo é a pedra principal da esquina.
\VS{21}Nele o edifício todo, bem ajustado, cresce para
{\ADD{ser}} um templo santo no Senhor.
\VS{22}Nele também vós sois juntamente edificados para
{\ADD{serdes}} morada de Deus em Espírito.

\par }\Chap{3}{\PP \VerseOne{1}Por isso eu, Paulo,
{\ADD{sou}} prisioneiro de Cristo Jesus, para o benefício de vós, gentios.
\VS{2}Se é que já ouvistes da responsabilidade
\FTNT{*}{{\FR 3:2 }
Ou, tradicionalmente: dispensação
} acerca da graça de Deus, que me foi dada para vós;
\VS{3}que por revelação me foi dado a conhecer este mistério (conforme já
{\ADD{vos}} escrevi um pouco;
\VS{4}quando o ledes, podeis entender o meu entendimento deste mistério de Cristo).
\VS{5}Esse mistério
\FTNT{†}{{\FR 3:5 }
“Esse mistério” foi posto explicitamente para simplificar a frase
} em outras gerações não foi dado a conhecer aos seres humanos
\FTNT{‡}{{\FR 3:5 }
Lit. filhos dos homens
} , como agora foi revelado pelo Espírito aos seus santos apóstolos e profetas:
\VS{6}{\ADD{isto é,
}} que os gentios são conjuntamente herdeiros,
{\ADD{membros}} de um mesmo corpo, e participantes da sua promessa em Cristo por meio do evangelho.
\VS{7}Eu fui feito servidor desse
{\ADD{evangelho}} conforme o dom da graça de Deus, que me foi dada segundo a operação do seu poder.
\VS{8}A mim, o menor de todos os santos, foi dada esta graça de anunciar entre os gentios, por meio do Evangelho, as inimagináveis riquezas de Cristo,
\VS{9}E para esclarecer a todos qual é a comunhão do mistério que desde as eras passadas esteve oculto em Deus, que criou todas as coisas por meio de Jesus Cristo,
\VS{10}para que a multiforme sabedoria de Deus seja agora manifestada pela igreja aos domínios e autoridades nos lugares celestes,
\VS{11}conforme o eterno propósito que ele cumpriu em Cristo Jesus, nosso Senhor,
\VS{12}no qual temos coragem
\FTNT{§}{{\FR 3:12 }
Ou: ousadia
} e acesso confiante pela fé nele.
\FTNT{*}{{\FR 3:12 }
Trad. alt.: pela fidelidade dele
}
\VS{13}Portanto, eu
{\ADD{vos}} peço que não vos desanimeis em minhas aflições por vós. Elas são a vossa glória.
\VS{14}Por causa disso me ponho de joelhos diante do Pai do nosso Senhor Jesus Cristo;
\VS{15}do qual toda a família nos céus e na terra recebe nome.
\VS{16}{\ADD{Oro}} para que, conforme a sua riquíssima glória,
\FTNT{†}{{\FR 3:16 }
Lit. as riquezas de sua glória
} ele vos conceda que sejais fortalecidos com poder pelo seu Espírito no interior de cada um,
\FTNT{‡}{{\FR 3:16 }
no interior de cada um
lit. no ser humano interior
}
\VS{17}para que Cristo habite em vossos corações pela fé.
{\ADD{Oro}} para que vós estejais enraizados e firmados no amor,
\VS{18}e assim possais compreender, com todos os santos, qual é a largura, comprimento, profundidade, e altura,
\VS{19}e conhecer o amor de Cristo, que excede o entendimento, para que sejais cheios de toda a plenitude de Deus.
\VS{20}E quanto àquele que é poderoso para fazer tudo muito mais abundantemente do que pedimos ou pensamos, segundo o poder que opera em nós,
\VS{21}a ele seja a glória na Igreja, em Cristo Jesus, por todas as gerações para todo o sempre, Amém!

\par }\Chap{4}{\PP \VerseOne{1}Portanto, eu, o prisioneiro no Senhor, rogo-vos que andeis como é digno do chamado com que fostes chamados,
\VS{2}com toda humildade e mansidão; com paciência, tolerando uns aos outros em amor.
\VS{3}Procurai guardar a unidade do Espírito pelo vínculo da paz.
\VS{4}Há um
{\ADD{só}} corpo, e um
{\ADD{só}} Espírito, como também fostes chamados em uma
{\ADD{só}} esperança do vosso chamado.
\VS{5}Um
{\ADD{só}} Senhor, uma
{\ADD{só}} fé, um
{\ADD{só}} batismo;
\VS{6}um
{\ADD{só}} Deus, e Pai de todos, que é acima de todos, por meio de todos, e em todos.
\VS{7}Mas a graça foi dada a cada um de nós conforme a medida do dom de Cristo.
\VS{8}Por isso diz: Quando ele subiu ao alto, levou cativo o cativeiro,
\FTNT{*}{{\FR 4:8 }
Ou: levou os cativos
} e deu dons às pessoas.
{\RQ{Salmos 68:18
}}
\VS{9}(O que significa este “ele subiu”, senão que ele também antes tinha descido às partes mais baixas da terra?)
\VS{10}Aquele que desceu é também o mesmo que subiu muito acima de todos os céus, para preencher
\FTNT{†}{{\FR 4:10 }
Trad. alt. cumprir
} todas as coisas.
\VS{11}E ele mesmo deu uns como apóstolos, outros como profetas, outros como evangelistas, e outros como pastores e instrutores.
\VS{12}para o aperfeiçoamento dos santos, para a obra do serviço,
\FTNT{‡}{{\FR 4:12 }
Ou, tradicionalmente: ministério
} para a edificação do corpo de Cristo;
\VS{13}até que todos cheguemos à unidade da fé e do conhecimento do Filho de Deus, à maturidade,
\FTNT{§}{{\FR 4:13 }
Lit. ao homem maduro
} à medida da estatura da plenitude de Cristo.
\FTNT{*}{{\FR 4:13 }
Ou: estatura completa de Cristo
}
\VS{14}O fim disso é que não mais sejamos crianças inconstantes, levadas de um lado para outro por todo vento de doutrina, pelo engano dos pessoas que, para enganar, usam de fraudes com astúcia.
\VS{15}Pelo contrário, sigamos a verdade em amor, e cresçamos em tudo naquele que é a cabeça: Cristo.
\VS{16}É dele que todo o corpo, bem ajustado e unido por todos os ligamentos, conforme a operação de cada parte na medida devida, proporciona o crescimento do corpo, para a sua própria edificação em amor.
\VS{17}Portanto digo isso, e dou testemunho no Senhor, que não andeis mais como os outros gentios andam, na futilidade das suas mentes.
\VS{18}Esses têm o entendimento nas trevas, e estão separados da vida de Deus pela ignorância que há neles, pela dureza do seu coração.
\VS{19}Como se tornaram insensíveis, eles se entregaram à sensualidade, para cometerem com avidez toda impureza.
\VS{20}Mas não foi assim que vós aprendestes de Cristo;
\VS{21}se é que o ouvistes, e fostes ensinados nele, como a verdade está em Jesus;
\VS{22}quanto ao comportamento passado, que deveis abandonar o velho ser humano, que se corrompe pelos maus desejos de engano;
\VS{23}que
{\ADD{deveis}} vos renovar no espírito da vossa mente;
\VS{24}e que
{\ADD{deveis}} vos revestir do novo ser humano, que é criado conforme Deus na verdadeira justiça e santidade.
\VS{25}Portanto, abandonai a mentira, e falai a verdade cada um ao seu próximo; pois somos membros uns dos outros.
\VS{26}Quando irardes, não pequeis;
\FTNT{†}{{\FR 4:26 }
Lit. Irai-vos, e não pequeis
} o sol não se ponha sobre a vossa ira;
\VS{27}nem deis lugar ao diabo.
\VS{28}O que furtava, não furte mais; em vez disso, trabalhe, operando com as mãos o bem, para que tenha o que repartir com o que tiver necessidade.
\VS{29}Não saia de vossa boca palavra imoral; mas sim a boa para a edificação conforme a necessidade; para que comunique graça aos que a ouvem.
\VS{30}E não entristeçais o Espírito Santo de Deus, no qual fostes selados para o dia da libertação.
\FTNT{‡}{{\FR 4:30 }
Ou: redenção
}
\VS{31}Toda amargura, raiva, ira, gritaria, e maledicência sejam tiradas do meio de vós, assim como toda malícia.
\VS{32}Em vez disso, sede benignos uns com os outros, misericordiosos, perdoando-vos uns aos outros, como também Deus vos perdoou em Cristo.

\par }\Chap{5}{\PP \VerseOne{1}Portanto, sede imitadores de Deus, como filhos amados;
\VS{2}e andai em amor, como também Cristo nos amou, e se entregou por nós como oferta e sacrifício de cheiro suave a Deus.
\VS{3}Mas o pecado sexual, e toda impureza ou cobiça, nem mesmo se nomeie entre vós, como convém a santos;
\VS{4}nem imoralidades, nem palavras tolas, nem escárnios, que não convêm; em vez disso, atos de gratidão.
\VS{5}Pois sabeis disto, que nenhum pecador no sexo, ou impuro, ou ganancioso, que é idólatra, tem herança no reino de Cristo e de Deus.
\VS{6}Ninguém vos engane com palavras vazias; pois por essas coisas a ira de Deus vem sobre os filhos da desobediência.
\VS{7}Portanto não sejais parceiros deles.
\FTNT{*}{{\FR 5:7 }
Ou: participantes com eles
}
\VS{8}Pois antes vós éreis trevas, mas agora
{\ADD{sois}} luz no Senhor; andai como filhos da luz
\VS{9}(porque o fruto do Espírito{\ADD{consiste}} em toda bondade, justiça, e verdade),
\VS{10}buscando descobrir o que agrada ao Senhor.
\VS{11}E não participeis das obras infrutíferas das trevas, pelo contrário, reprovai-as.
\VS{12}pois é vergonhoso até dizer o que eles fazem em segredo.
\VS{13}Mas todas essas coisas, quando são reprovadas, tornam-se visíveis pela luz, porque tudo o que se se torna visível é luz.
\VS{14}Por isso se diz: Desperta tu, que dormes, e levanta-te dos mortos, e Cristo te iluminará.
\VS{15}Portanto vede cuidadosamente como vos conduzis,
\FTNT{†}{{\FR 5:15 }
Lit. andais
} não como tolos, mas sim como sábios;
\VS{16}aproveitai o tempo, porque os dias são maus.
\VS{17}Por isso, não sejais imprudentes, mas entendei qual é a vontade do Senhor.
\VS{18}E não fiqueis bêbados com vinho, em que há devassidão, mas enchei-vos do Espírito;
\VS{19}falando entre vós com salmos, hinos e cânticos espirituais; cantando e louvando ao Senhor no vosso coração;
\VS{20}agradecendo sempre por tudo a Deus, o Pai, no nome do nosso Senhor Jesus Cristo;
\VS{21}sujeitando-vos uns aos outros no temor a Deus.
\VS{22}Mulheres, sujeitai-vos aos vossos próprios maridos, assim como ao Senhor;
\VS{23}porque o marido é o cabeça da mulher, como também Cristo é o cabeça da Igreja, e ele é o salvador do corpo.
\VS{24}Mas, assim como a Igreja está sujeita a Cristo, assim também as mulheres
{\ADD{estejam}} em tudo sujeitas aos seus próprios maridos.
\VS{25}Maridos, amai as vossas próprias esposas, assim como também Cristo amou a Igreja, e se entregou por ela;
\VS{26}a fim de a santificar, havendo
{\ADD{a}} purificado com a lavagem com água, pela palavra,
\VS{27}a fim de apresentar a si mesmo uma igreja gloriosa, sem mancha, nem ruga, nem algo semelhante; mas sim, santa e irrepreensível.
\VS{28}Assim os maridos devem amar as suas próprias esposas como os seus próprios corpos. Quem ama a sua esposa, ama a si mesmo;
\VS{29}pois ninguém jamais odiou a sua própria carne; mas a alimenta e sustenta, assim como também o Senhor à Igreja.
\VS{30}pois somos membros de seu corpo, da sua carne, e dos seus ossos.
\VS{31}Por isso o homem deixará oseu pai e
{\ADD{a sua}} mãe, e se ajuntará com a sua mulher; e os dois serão uma
{\ADD{só}} carne.
\VS{32}Esse é um grande mistério, mas estou dizendo quanto a Cristo e à Igreja.
\VS{33}Assim também vós, cada um individualmente, ame a sua própria esposa como a si mesmo, e a mulher respeite o marido.

\par }\Chap{6}{\PP \VerseOne{1}Filhos, sede obedientes aos vossos pais no Senhor, porque isso é justo.
\VS{2}Honra teu pai, e
{\ADD{tua}} mãe (que é o primeiro mandamento com promessa);
\VS{3}para que te vá bem, e vivas muito tempo sobre a terra.
{\RQ{Deuteronômio 5:16
}}
\VS{4}Pais, não provoqueis à ira os vossos filhos, mas criai-os na disciplina e correção do Senhor.
\VS{5}Servos, obedecei aos
{\ADD{vossos}} senhores segundo a carne, com temor e tremor, em sinceridade do vosso coração, assim como a Cristo;
\VS{6}não servindo somente quando vistos, como para agradar pessoas; mas sim, como servos de Cristo, fazendo de coração a vontade de Deus;
\VS{7}servindo de boa vontade como ao Senhor, e não aos pessoas.
\VS{8}Pois sabeis que cada um receberá do Senhor todo o bem que fizer, seja servo, seja livre.
\VS{9}Senhores, fazei o mesmo a eles, abandonando as ameaças, pois sabeis também que o Senhor vosso e deles está nos céus, e
{\ADD{que}} com ele não há acepção de pessoas.
\VS{10}Por fim, meus irmãos, fortalecei-vos no Senhor e na força do seu poder.
\VS{11}Revesti-vos de toda a armadura de Deus, para que possais estar firmes contra as ciladas do diabo.
\VS{12}Pois temos de lutar não contra carne e sangue, mas sim contra os domínios e poderes, contra os senhores das trevas deste mundo, contra os males espirituais nos lugares celestes.
\VS{13}Portanto pegai toda a armadura de Deus, para que possais resistir no dia mal, e depois que fizerdes tudo, continuar firmes.
\VS{14}Estai, pois, firmes, com a vossa cintura envolvida com o cinturão da verdade, e vestidos com a couraça da justiça;
\VS{15}e calçados os pés com a prontidão do evangelho da paz.
\VS{16}Sobretudo pegai o escudo da fé, com o qual podereis apagar todas as flechas inflamadas do maligno.
\VS{17}Pegai também o capacete da salvação, e a espada do Espírito, que é a palavra de Deus;
\VS{18}orando em todo tempo com toda oração e súplica no Espírito, e vigiando nisso com toda perseverança e súplica por todos os santos,
\VS{19}e por mim; a fim de que, quando eu abrir a boca, me seja dada palavra com ousadia, para tornar conhecido o mistério do evangelho,
\VS{20}pelo qual sou embaixador acorrentado;
\FTNT{*}{{\FR 6:20 }
Lit. em correntes
} para que dele eu possa falar com ousadia, como devo falar.
\VS{21}Para que também vós possais saber o que acontece comigo,
{\ADD{e}} o que eu faço, Tíquico, o irmão amado, e fiel servidor do Senhor vos informará de tudo.
\VS{22}Eu o enviei até vós para isso, para que saibais o que acontece conosco, e ele conforte os vossos corações.
\VS{23}Paz
{\ADD{seja}} com os irmãos, e amor com fé, da parte de Deus Pai, e do Senhor Jesus Cristo.
\VS{24}A graça
{\ADD{seja}} com todos os que amam ao nosso Senhor Jesus Cristo com amor que não falha.
\FTNT{†}{{\FR 6:24 }
com amor que não falha
Lit. em incorrupção
} Amém!{\ADD{Escrita de Roma para os efésios, e enviada por Tíquico}}
\par }